\section{Zone Architecture}\label{sec:architecture}
\subsection{Input Scope and Instruction Representation}
Zone's payload implementation is written in Python and uses Capstone
for disassembly and Keystone for assembly. The project is publicly
available.\footnote{\url{https://github.com/jakydibe/Zone}.}
Its x64 and x86 paths share the same broad organization but require
architecture-specific handling of operands and references. Support in
one path is therefore not evidence of support in the other.

Instruction objects retain an association with the original decoding
and the current representation. Address-indexed lookup and tracked
control-flow references allow layout updates to be propagated through
the instruction sequence. A rewriting operation can change the size of
an instruction, its address, or the physical relation between blocks;
the final output must reflect the resulting layout rather than the
addresses in the first decoding.

The initial decoding is itself fallible. Decoder options that skip
undecodable bytes allow processing to continue, but do not establish
which successfully decoded bytes are genuine code. Full decoding
coverage is a useful structural check and remains insufficient as a
semantic criterion. The extended data layer addresses a different part
of this problem by retaining explicit candidate data and reference
information.

\subsection{The Published PII Design}\label{sec:pii}
PII was introduced in the conference article~\cite{zone_secrypt_2026}.
It separates logical execution order from the physical placement of
instruction units. The published representation uses wrapped original
operations, a redirector, and an offset table. The table identifies the
next logical unit after the current operation; physical adjacency is
therefore not the sole representation of succession.
Figure~\ref{fig:pii_layout} retains the original layout as architectural
context. It describes the conference design rather than a new method
introduced in this extension.
\begin{figure}[t]
\centering
\begin{tikzpicture}[
    font=\scriptsize,
    node distance=0.0cm,
    every node/.style={
        rotate=90,
        minimum width=3.3cm,
        minimum height=0.7cm,
        text centered,
        draw,
        fill=gray!10,
        rounded corners,
        inner xsep=2pt,
        inner ysep=6pt
    }
]

\node (stubCall)  [fill=blue!10]    {Call Stub 1};
\node (payload)   [below=of stubCall, fill=green!10]  {Wrapped Inst. / payload};
\node (pad1)      [below=of payload]                  {PAD (NOP)};
\node (stub2)     [below=of pad1, fill=yellow!20]     {Stub 2};
\node (pad2)      [below=of stub2]                    {PAD (NOP)};
\node (stub1)     [below=of pad2, fill=blue!10]       {Stub 1};
\node (pad3)      [below=of stub1]                    {PAD (NOP)};
\node (payloadPtr)[below=of pad3, fill=red!20]        {Pointer to Payload};
\node (pad4)      [below=of payloadPtr]               {PAD (NOP)};
\node (memZone)   [below=of pad4, fill=purple!10]     {Memory Zone};

\end{tikzpicture}
\Description{A schematic separates the initial call, instruction wrappers, redirector components, padding, offset information, and memory region.}
\caption{PII payload layout. Adapted from Fig.~1 of~\cite{zone_secrypt_2026}.}
\label{fig:pii_layout}
\end{figure}

The state used by dispatch must be preserved independently of the state
used by the original operation. Control-flow instructions and
architecture-specific addressing also require dedicated handling.
The extension distinguishes the original PII design from later
implementation revisions. It also distinguishes a dispatcher-based
representation from single-instruction permutation with explicit
successors. These paths are not interchangeable experimental labels.

\subsection{Extended Architecture}
Figure~\ref{fig:extended_architecture} organizes the added functionality
into identification, deterministic data handling, transformation, and
validation. The identification methods are alternatives. Both supply
candidate regions to the data layer, while the LLM transformation path
uses an assembly representation prepared by the engine. The model is
not responsible for the entire binary-writing infrastructure.

\begin{figure}[t]
\centering
\begin{tikzpicture}[font=\scriptsize,
 box/.style={draw,rounded corners,align=center,minimum height=8mm,text width=30mm},
 arr/.style={-{Stealth[length=2mm]},thick}]
\node[box,fill=gray!10] (input) at (0,0) {Raw payload\\and architecture};
\node[box,fill=blue!7] (sym) at (-2.1,-1.4) {Symbolic and heuristic\\identification};
\node[box,fill=blue!7] (llmdata) at (2.1,-1.4) {LLM-based\\identification};
\node[box,fill=green!7,text width=42mm] (data) at (0,-2.9) {Shared deterministic\\data and reference handling};
\node[box,fill=orange!8] (native) at (-2.1,-4.4) {Native transformations\\and wrapping};
\node[box,fill=orange!8] (llm) at (2.1,-4.4) {LLM-assisted\\assembly transformation};
\node[box,fill=gray!10,text width=42mm] (validation) at (0,-5.9) {Structural validation\\and functional assessment};
\draw[arr] (input)--(sym); \draw[arr] (input)--(llmdata);
\draw[arr] (sym)--(data); \draw[arr] (llmdata)--(data);
\draw[arr] (data)--(native); \draw[arr] (data)--(llm);
\draw[arr] (native)--(validation); \draw[arr] (llm)--(validation);
\end{tikzpicture}
\caption{Conceptual organization of the extended system. Identification
and transformation paths are selected according to the experimental
condition; the diagram does not assert that every combination is supported.}
\label{fig:extended_architecture}
\Description{Two alternative identification paths feed shared data handling,
followed by native or LLM-assisted transformation and validation.}
\end{figure}

The implementation distinguishes original input identity, intermediate
code and data representations, and final output identity. This matters
for both reference handling and measurement. For example, similarity
against a canonicalized intermediate representation measures a different
change from similarity against the original file. Likewise, a data
analysis retrieved from a cache is not a new symbolic exploration.

\begin{table}[t]
\caption{Implementation scope of the extensions in this working draft.
Availability does not imply completed empirical validation.}
\label{tab:scope}
\centering\small
\begin{tabularx}{\linewidth}{@{}Yll@{}}
\toprule
Component & x64 & x86 \\
\midrule
Symbolic/heuristic identification & Implemented & Implemented \\
LLM data identification & Implemented & Implemented \\
LLM EQ/NOP/BCF/IBP & Implemented & Implemented \\
Assembly wrapping & Implemented & \dev \\
Extended quantitative assessment & In progress & In progress \\
\bottomrule
\end{tabularx}
\end{table}

Table~\ref{tab:scope} records the current development boundary.
The support matrix will be finalized against the version used for the
reported campaigns. Ongoing development, especially the x86 wrapping
port, is not merged into an already completed experimental condition.
